\section{4 Founder Profiles, 4 Opportunity Maps}

\subsection{Beyond the Monoculture Playbook}

A theory of entrepreneurship that works only for software engineers is not a general theory. Consider 4 distinct starting archetypes:

\subsection{Archetype Profiles}

\textbf{1. The Technical Specialist}
\begin{itemize}[leftmargin=11pt,itemsep=1pt,topsep=1pt,parsep=0pt]
  \item \textit{Assets:} Strong engineering ability, deep insight into internal software bottlenecks, small developer circle, minimal cash.
  \item \textit{Weakness:} Few commercial contacts, zero enterprise sales experience.
  \item \textit{Opening:} Developer tools, narrow technical diagnostics, productized workflow scripts.
  \item \textit{False inference:} ``Because I can build the system, commercial demand must exist.''
  \item \textit{Best early proof:} Working prototype solving an acute pain; paid pilot agreements.
  \item \textit{1st bottleneck:} Reaching the budget holder who can sign a check.
\end{itemize}

\textbf{2. The Experienced Service Professional}
\begin{itemize}[leftmargin=11pt,itemsep=1pt,topsep=1pt,parsep=0pt]
  \item \textit{Assets:} Deep workflow knowledge, established trust with clients, diagnostic instincts.
  \item \textit{Weakness:} Limited coding skill, high personal delivery reliance.
  \item \textit{Opening:} Specialized consulting, regulatory compliance audits, productized services.
  \item \textit{False inference:} ``I must raise venture capital and build an app before I can start.''
  \item \textit{Best early proof:} High-margin fixed-scope projects with repeatable diagnostics.
  \item \textit{1st bottleneck:} Decoupling founder labor from value delivery to preserve margins.
\end{itemize}

\textbf{3. The Consumer-Brand Builder}
\begin{itemize}[leftmargin=11pt,itemsep=1pt,topsep=1pt,parsep=0pt]
  \item \textit{Assets:} Exceptional aesthetic taste, subculture credibility, supplier access.
  \item \textit{Weakness:} Thin working capital, little supply chain operations experience.
  \item \textit{Opening:} Focused SKU line, small-batch pre-orders, boutique channel placement.
  \item \textit{False inference:} ``Refined aesthetics alone guarantee a premium margin.''
  \item \textit{Best early proof:} Repeat purchases, viable landed unit costs, low return rates.
  \item \textit{1st bottleneck:} Inventory financing and defect control under scale.
\end{itemize}

\textbf{4. The Scientist or Researcher}
\begin{itemize}[leftmargin=11pt,itemsep=1pt,topsep=1pt,parsep=0pt]
  \item \textit{Assets:} Rare scientific expertise, publications, peer network, proprietary IP.
  \item \textit{Weakness:} IP ownership encumbrances, long research cycles, buyer unfamiliarity.
  \item \textit{Opening:} Research instrumentation, contract R\&D, patent licensing.
  \item \textit{False inference:} ``Breakthrough scientific novelty automatically commands a market.''
  \item \textit{Best early proof:} Reproducible results, freedom-to-operate clearance, pilot buyer.
  \item \textit{1st bottleneck:} Capital intensity and commercial translation risk.
\end{itemize}

\textbf{Founders with different starting configurations must not begin with the same playbook: their best opening experiments differ because their starting assets and binding bottlenecks differ.}
